\chapter*{符号说明}
\markboth{\sffamily\slshape 符号说明}{\sffamily\slshape 符号说明}

%% Consistent section title command for this page
\newcommand{\notationsection}[1]{%
  \vspace{1.2em}%
  \noindent{\textbfs{#1}}%
  \vspace{0.4em}%
  \par\nobreak
}

%% Shared table setup
\newcommand{\notationtable}[1]{%
  \bgroup
  \def\arraystretch{1.25}%
  \noindent\begin{tabular}{@{}p{0.3\linewidth}@{\hspace{1.5em}}p{0.64\linewidth}@{}}
  #1
  \end{tabular}
  \egroup
}

%%% ---------------------------------------------------------------
\notationsection{数与数组}
\notationtable{
$a$                                       & 标量。\\
$\rva$                                    & 列向量（例如 $\rva\in\mathbb{R}^D$）。\\
$\rmA$                                    & 矩阵（例如 $\rmA\in\mathbb{R}^{m\times n}$）。\\
$\rmA^\top$                               & $\rmA$ 的转置。\\
$\Tr(\rmA)$                               & $\rmA$ 的迹。\\
$\rmI_D$                                  & $D\times D$ 的单位矩阵。\\
$\rmI$                                    & 单位矩阵；维数由上下文确定。\\
$\mathrm{diag}(\rva)$                     & 以 $\rva$ 为对角元的对角矩阵。\\
$\bm{\phi},\,\bm{\theta}$                & 可学习的神经网络参数。\\
$\bm{\phi}^{\times},\,\bm{\theta}^{\times}$ & 训练后的参数（在推理阶段固定）。\\
$\bm{\phi}^{*},\,\bm{\theta}^{*}$        & 优化问题的最优参数。\\
}

%%% ---------------------------------------------------------------
\notationsection{微积分}
\notationtable{
$\dfrac{\partial \rvy}{\partial \rvx}$
    & $\rvy$ 对 $\rvx$ 的偏导数（逐分量）。\\[6pt]
$\dfrac{\diff \rvy}{\diff \rvx}$ \text{或} $\mathrm{D}\rvy(\rvx)$
    & $\rvy$ 对 $\rvx$ 的（Fréchet）全导数。\\[6pt]
$\nabla_\rvx y$
    & 标量函数 $y:\mathbb{R}^D \to\mathbb{R}$ 的梯度；$\mathbb{R}^D$ 中的列向量。\\
$\dfrac{\partial \rmF}{\partial \rvx}$ \text{或} $\nabla_\rvx \rmF$
    & $\rmF:\mathbb{R}^n \to\mathbb{R}^m$ 的雅可比矩阵；形状为 $m\times n$。\\[6pt]
$\nabla\cdot\rvy$
    & 向量场 $\rvy:\mathbb{R}^D \to\mathbb{R}^D$ 的散度；一个标量。\\
$\nabla^2_\rvx f(\rvx)$ \text{或} $\rmH(f)(\rvx)$
    & $f:\mathbb{R}^D \to\mathbb{R}$ 的海森矩阵；形状为 $D\times D$。\\
$\displaystyle\int f(\rvx)\diff\rvx$
    & $f$ 在 $\rvx$ 的定义域上的积分。\\
}

%%% ---------------------------------------------------------------
\notationsection{概率论与信息论}
\notationtable{
$p(\rvx)$                   & 连续向量 $\rvx$ 上的密度/分布。\\
$\Pr(\rvx = \rva)$，或非正式地写作 $\Pr(\rvx)$
    & 离散情形下的点概率。\\
$\Pr(\rvx = \rva |\rvy = \rvb)$，或非正式地写作 $\Pr(\rvx |\rvy)$
    & 离散情形下的条件点概率。\\
$p_{\mathrm{data}}$        & 数据分布。\\
$p_{\mathrm{prior}}$       & 先验分布（例如标准正态分布）。\\
$p_{\mathrm{src}}$         & 源分布。\\
$p_{\mathrm{tgt}}$         & 目标分布。\\
$\rva \sim p$              & 随机向量 $\rva$ 服从分布 $p$。\\
$\E_{\rvx\sim p}\big[\rvf(\rvx)\big]$
    & $\rvf(\rvx)$ 在 $p(\rvx)$ 下的期望。\\[3pt]
\parbox[t]{0.22\linewidth}{\raggedright
  $\E\big[\rvf(\rvx)\,|\,\rvz\big]$，\\[2pt]
  或 $\E_{\rvx\sim p(\cdot|\rvz)}\big[\rvf(\rvx)\big]$}
    & 给定 $\rvz$ 时 $\rvf(\rvx)$ 的条件期望，其中 $\rvx$ 服从 $p(\cdot|\rvz)$。\\[3pt]
$\Var\big(\rvf(\rvx)\big)$
    & 在 $p(\rvx)$ 下的方差。\\
$\Cov\big(\rvf(\rvx),\rvg(\rvx)\big)$
    & 在 $p(\rvx)$ 下的协方差。\\
$\mathcal{D}_{\mathrm{KL}}\left(p\,\Vert\, q\right)$
    & 从 $q$ 到 $p$ 的 Kullback--Leibler 散度。\\
$\bm{\epsilon}\sim\mathcal{N}(\bm{0},\rmI)$
    & 标准正态分布的样本。\\
$\mathcal{N}(\rvx;\vmu,\mSigma)$
    & $\rvx$ 上的高斯分布，均值为 $\vmu$，协方差为 $\mSigma$。\\
}

%%% ---------------------------------------------------------------
\vspace{1.2em}

\paragraph{记号约定。}
本书通篇使用若干约定，在此一并说明。

\subparagraph{I. 随机变量及其取值。}
我们对随机向量及其取值（实现值）使用相同的符号。
这一约定在深度学习和生成建模中很常见，可以使记号更紧凑。其具体含义由上下文确定。

对于密度函数，$p(\rvx)$ 表示分布的函数形式，而不是在某个具体样本处的取值。当
需要在某个具体点处求值时，我们会明确说明。
条件密度由上下文来理解：在 $p(\rvx|\rvy)$ 中，固定
$\rvy$ 得到关于 $\rvx$ 的密度函数，而固定 $\rvx$ 则得到关于
$\rvy$ 的函数。

\subparagraph{II. 条件期望。}
表达式 $\E[\rvf(\rvx)|\rvz]$ 表示关于 $\rvz$ 的函数，
给出在给定 $\rvz$ 的条件下 $\rvf(\rvx)$ 的期望值。
当以某个具体的取值为条件时，我们记作
$\E[\rvf(\rvx)|\rmZ = \rvz]$。等价地，这可以表示为
关于条件分布的积分：
\[
\E_{\rvx \sim p(\cdot|\rvz)}[\rvf(\rvx)]
= \int \rvf(\rvx)\, p(\rvx|\rvz)\, \diff\rvx.
\]
这种区分可以澄清 $\rvz$ 是被当作定义函数 $\rvz \mapsto \E[\rvf(\rvx)|\rvz]$ 的变量，
还是被当作对该函数求值的固定点。

\subparagraph{III. 隐含的时间下标。}
由于本书中的许多对象都以时间作为下标，把每个
时间变量都显式写出会使公式显得繁重。当相关的时间
下标能从上下文中清楚看出时，我们将其省略。例如，
\[
\E[\rvx_s | \rvx_t]
\quad\text{简写为}\quad
\E[\rvx_s | \rvx_t,\, s,\, t],
\]
即在时间 $t$ 给定 $\rvx_t$ 的条件下，$\rvx_s$ 在时间 $s$ 处的
平均，其中 $s$ 和 $t$ 的含义由上下文确定。
类似地，$p(\rvx_s | \rvx_t)$ 表示给定 $\rvx_t$ 时 $\rvx_s$ 的条件分布，
两个时间下标都隐含其中。

\subparagraph{IV. 高斯扰动与依分布相等。}
对于固定的下标 $t$，我们经常使用如下高斯扰动核：
\[
p_t(\rvx_t|\rvx_0)
= \mathcal{N}\big(\rvx_t;\, \alpha_t \rvx_0,\, \sigma_t^2 \rmI\big),
\]
其中 $\rvx_0 \sim p_{\mathrm{data}}$，$\alpha_t$ 和 $\sigma_t$ 是
确定性的标量。我们等价地将其写成\emph{依分布成立}的等式
\[
\rvx_t \stackrel{d}{=} \alpha_t \rvx_0 + \sigma_t \beps,
\qquad
\beps \sim \mathcal{N}(\bm{0}, \rmI),
\]
其中 $\beps$ 与 $\rvx_0$ 独立。符号
``$\stackrel{d}{=}$'' 表示两边具有相同的概率律，因此
\[
    \E[\phi(\rvx_t)] = \E[\phi(\alpha_t \rvx_0 + \sigma_t \beps)]
\]
对于任意测试函数 $\phi$ 均成立。为简便起见，我们将直接
写作 $\rvx_t = \alpha_t \rvx_0 + \sigma_t \beps$，根据上下文理解为
依分布相等或一次样本实现；本书通篇采用这一简写。
